% ------------------------------------------------------------
% 2. Metodología
% ------------------------------------------------------------

\subsection{Datos}
\label{sec:datos}

Los datos provienen de la Tabla~6 de \citet{Lakdawala2025} y corresponden a las 80 corridas descritas en la introducción. La variable respuesta es la eficiencia hidráulica, expresada como proporción (entre 0 y 1). La Tabla~\ref{tab:variables} resume las variables y el nombre con el que aparecen en la base de datos.

\begin{table}[ht]
\centering
\caption{Variables del estudio.}
\label{tab:variables}
\begin{tabular}{llll}
\toprule
Variable & Nombre en R & Tipo & Niveles \\
\midrule
Eficiencia hidráulica ($y$)            & \texttt{eficiencia} & Respuesta numérica   & 0.6025 a 0.7165 \\
Ángulo de la boquilla ($\alpha$)       & \texttt{angulo}     & Numérica (grados)    & 90, 92.5, 95 \\
Diámetro de la boquilla ($d_j$)        & \texttt{diametro}   & Numérica (mm)        & 10, 12, 14 \\
Distancia boquilla-rodete ($N_{sod}$)  & \texttt{distancia}  & Numérica (mm)        & 40, 50, 60 \\
Tipo de rodete                         & \texttt{rodete}     & Categórica           & Pelton (PTR), Turgo (TTR) \\
Número de boquillas ($n$)              & \texttt{boquillas}  & Categórica           & 1, 2 \\
\bottomrule
\end{tabular}
\end{table}

Al revisar la Tabla~6 se encontró que en dos corridas los niveles publicados no coinciden con los valores predichos que reportan los mismos autores: en la corrida 12 el ángulo corresponde a 95° (no a 90°) y en la corrida 15 la distancia corresponde a 60~mm (no a 40~mm). En este trabajo se usan esos dos niveles corregidos; con ellos el diseño queda completo y el modelo reproduce exactamente los resultados del artículo.

Los datos y el código del análisis están disponibles en \url{https://github.com/mariasaurius/Articulo_regresion}. El script \texttt{anal\_descrip.R} produce las tablas y figuras del análisis descriptivo, y el script \texttt{mod\_lin.R} las del modelo de regresión. Ambos leen los datos desde el repositorio y aplican la misma corrección.

\subsection{Modelo de regresión}
\label{sec:modelo}

El objetivo es encontrar una ecuación que prediga la eficiencia a partir de las cinco variables. El análisis descriptivo de la Sección~\ref{sec:descriptivo} mostró dos cosas que esa ecuación debe poder representar.

\begin{itemize}
  \item \textbf{Curvatura.} La eficiencia no baja en línea recta cuando aumenta el ángulo. Según la Tabla~\ref{tab:niveles}, la eficiencia media pasa de 0.6813 a 90° a 0.6478 a 92.5° (baja 0.0335) y luego a 0.6389 a 95° (baja solo 0.0089). Los dos pasos son de 2.5°, pero la primera caída es casi cuatro veces la segunda. Si la relación fuera una recta, las dos caídas serían iguales. Con la distancia pasa algo parecido.
  \item \textbf{Interacciones.} El efecto de una variable parece depender de otra; por ejemplo, el diámetro casi no cambia la eficiencia con 90° pero sí con 95° (Tabla~\ref{tab:angulo_diametro}).
\end{itemize}

Un modelo con solo rectas no puede representar ninguna de las dos cosas. Por eso, siguiendo a \citet{Lakdawala2025}, se usa un modelo de segundo orden, que es el modelo típico de la metodología de superficie de respuesta.

\subsubsection{Codificación de las variables}

Antes de ajustar el modelo, las variables se pasan a una escala común. Cada variable numérica se codifica de modo que su nivel bajo valga $-1$, el medio $0$ y el alto $1$, y las dos categóricas se codifican con $+1$ y $-1$. La Tabla~\ref{tab:codificacion} resume las cinco variables codificadas, que se llaman $A$, $B$, $C$, $D$ y $E$, como en el artículo original.

\begin{table}[ht]
\centering
\caption{Codificación de las variables del modelo.}
\label{tab:codificacion}
\small
\begin{tabular}{clllc}
\toprule
Símbolo & Variable & Niveles originales & Fórmula & Codificada \\
\midrule
$A$ & Ángulo de la boquilla     & 90°, 92.5°, 95°    & $(\text{ángulo} - 92.5)/2.5$    & $-1$, $0$, $1$ \\
$B$ & Diámetro de la boquilla   & 10, 12, 14 mm      & $(\text{diámetro} - 12)/2$      & $-1$, $0$, $1$ \\
$C$ & Distancia boquilla-rodete & 40, 50, 60 mm      & $(\text{distancia} - 50)/10$    & $-1$, $0$, $1$ \\
$D$ & Tipo de rodete            & Pelton, Turgo      & Pelton $= +1$, Turgo $= -1$     & $+1$, $-1$ \\
$E$ & Número de boquillas       & 1, 2               & Una $= +1$, dos $= -1$          & $+1$, $-1$ \\
\bottomrule
\end{tabular}
\end{table}

La fórmula de las numéricas resta el valor central y divide por la mitad del rango. Por ejemplo, para el ángulo: $(90 - 92.5)/2.5 = -1$, $(92.5 - 92.5)/2.5 = 0$ y $(95 - 92.5)/2.5 = 1$. La codificación tiene tres ventajas.

\begin{enumerate}
  \item \textbf{Permite comparar el tamaño de los coeficientes.} En las unidades originales, el coeficiente de una variable dice cuánto cambia la eficiencia por cada unidad de esa variable: por cada grado de ángulo o por cada milímetro de distancia. Pero esas unidades no son comparables entre sí: un grado es la quinta parte del rango del ángulo (de 90° a 95°), mientras que un milímetro es apenas la vigésima parte del rango de la distancia (de 40 a 60~mm). Una variable podría tener un coeficiente pequeño solo porque su unidad es pequeña, no porque importe poco. Con la codificación, una unidad es lo mismo para todas las variables: ir del centro al extremo de su rango. Así, el coeficiente más grande en valor absoluto corresponde a la variable que más mueve la eficiencia dentro del experimento.
  \item \textbf{Le da sentido al intercepto.} El intercepto es la predicción cuando todas las variables valen cero. En las unidades originales eso sería un ángulo de 0° y una boquilla de 0~mm, condiciones que no tienen sentido. Con la codificación, $A = B = C = 0$ es el centro del diseño (92.5°, 12~mm y 50~mm), y $D = E = 0$ es el punto medio entre $+1$ y $-1$, es decir, el promedio entre los dos rodetes y entre una y dos boquillas. Por eso el intercepto es la eficiencia predicha en el centro del diseño, promediando rodetes y número de boquillas.
  \item \textbf{Permite comparar con el artículo.} Es la escala que usa el programa Design Expert, con el que los autores ajustaron su modelo.
\end{enumerate}

\subsubsection{El modelo de segundo orden}

La forma general de un modelo de segundo orden con $k$ variables es
\begin{equation}
y = \beta_0 + \sum_{i=1}^{k} \beta_i x_i + \sum_{i=1}^{k} \beta_{ii} x_i^2 + \sum_{i<j} \beta_{ij} x_i x_j + \varepsilon,
\label{eq:general}
\end{equation}
es decir, tiene un intercepto, un término de primer orden para cada variable ($x_i$), un término cuadrático para cada variable ($x_i^2$) y un término de interacción para cada pareja de variables ($x_i x_j$) \citep{Montgomery2005}. Con $k$ variables hay $k$ términos de primer orden, $k$ cuadráticos y $k(k-1)/2$ interacciones; con $k = 5$ serían $1 + 5 + 5 + 10 = 21$ coeficientes.

El modelo de este trabajo tiene 19 y no 21 porque le faltan los cuadrados de $D$ y $E$. Como esas variables solo valen $+1$ o $-1$, su cuadrado siempre es $(+1)^2 = (-1)^2 = 1$: una columna de unos, igual a la del intercepto, que no aporta información nueva. Por eso solo las variables numéricas, que tienen tres niveles, llevan término cuadrático. El modelo queda así:
\begin{align}
y = \beta_0
&+ \underbrace{\beta_1 A + \beta_2 B + \beta_3 C + \beta_4 D + \beta_5 E}_{\text{5 efectos principales}}
 + \underbrace{\beta_{11} A^2 + \beta_{22} B^2 + \beta_{33} C^2}_{\text{3 términos cuadráticos}} \nonumber \\
&+ \underbrace{\beta_{12} AB + \beta_{13} AC + \cdots + \beta_{45} DE}_{\text{10 interacciones dobles}}
 + \varepsilon, \qquad \varepsilon \sim N(0, \sigma^2),
\label{eq:modelo2}
\end{align}
con $1 + 5 + 3 + 10 = 19$ coeficientes. Cada grupo de términos cumple una función:
\begin{itemize}
  \item \textbf{Intercepto} ($\beta_0$): la eficiencia predicha en el centro del diseño, como se explicó arriba.
  \item \textbf{Efectos principales} ($\beta_1, \ldots, \beta_5$): el efecto de cada variable por sí sola, con las demás en el centro. Como cada variable va de $-1$ a $1$, su coeficiente es la mitad de la diferencia entre la eficiencia predicha en el nivel alto y en el nivel bajo. Por ejemplo, $\beta_1$ es la mitad de la diferencia entre la eficiencia predicha con 95° y con 90°. Si es negativo, la eficiencia es menor en el nivel alto que en el bajo.
  \item \textbf{Términos cuadráticos} ($\beta_{11}, \beta_{22}, \beta_{33}$): permiten que la relación sea una curva y no una recta, es decir, que la caída de $-1$ a $0$ sea distinta de la caída de $0$ a $1$, como pasa con el ángulo.
  \item \textbf{Interacciones dobles} ($\beta_{12}, \ldots, \beta_{45}$): el producto de dos variables permite que el efecto de una dependa del valor de la otra. Con cinco variables hay $5 \times 4 / 2 = 10$ parejas posibles: $AB$, $AC$, $AD$, $AE$, $BC$, $BD$, $BE$, $CD$, $CE$ y $DE$.
\end{itemize}

Aunque la ecuación tiene cuadrados y productos, sigue siendo un modelo de \emph{regresión lineal}. La palabra \emph{lineal} se refiere a los coeficientes $\beta$, no a las variables: la eficiencia es una suma de coeficientes multiplicados por columnas de datos. Una columna como $A^2$ o $AB$ se calcula a partir de los datos antes de ajustar el modelo y, para el método de estimación, es una variable más. Un modelo no lineal sería, por ejemplo, $y = \beta_0 e^{\beta_1 A}$, donde un coeficiente aparece dentro de una exponencial.

\subsubsection{La matriz de diseño}

Escrito en forma matricial, el modelo es $\mathbf{y} = \mathbf{X}\bm{\beta} + \bm{\varepsilon}$, donde $\mathbf{y}$ tiene las 80 eficiencias, $\bm{\beta}$ los 19 coeficientes y $\mathbf{X}$ es la matriz de diseño, con una fila por corrida y una columna por coeficiente. La Tabla~\ref{tab:matrizX} muestra cuatro filas de $\mathbf{X}$ (sección~5.1 de \texttt{mod\_lin.R}). La primera columna es de unos (para el intercepto); las cinco siguientes son las variables codificadas, y las demás se obtienen elevando al cuadrado o multiplicando esas columnas. Por ejemplo, en la corrida~2 ($A = -1$, $B = 1$) la columna $AB$ vale $(-1)(1) = -1$, y en la corrida~5, que está en el centro del diseño, todas las columnas que contienen $A$, $B$ o $C$ valen cero.

\begin{table}[ht]
\centering
\caption{Cuatro filas de la matriz de diseño $\mathbf{X}$: corrida 1 (95°, 12~mm, 50~mm, Pelton, dos boquillas), corrida 2 (90°, 14~mm, 40~mm, Turgo, una boquilla), corrida 5 (punto central, Turgo, dos boquillas) y corrida 11 (90°, 10~mm, 40~mm, Turgo, dos boquillas).}
\label{tab:matrizX}
\setlength{\tabcolsep}{2.6pt}
\footnotesize
\begin{tabular}{c|r|rrrrr|rrr|rrrrrrrrrr}
\toprule
Corr. & 1 & $A$ & $B$ & $C$ & $D$ & $E$ & $A^2$ & $B^2$ & $C^2$ & $AB$ & $AC$ & $AD$ & $AE$ & $BC$ & $BD$ & $BE$ & $CD$ & $CE$ & $DE$ \\
\midrule
1  & 1 &  1 &  0 &  0 &  1 & $-1$ & 1 & 0 & 0 &  0 & 0 & 1 & $-1$ &  0 &  0 & 0 & 0 &  0 & $-1$ \\
2  & 1 & $-1$ &  1 & $-1$ & $-1$ &  1 & 1 & 1 & 1 & $-1$ & 1 & 1 & $-1$ & $-1$ & $-1$ & 1 & 1 & $-1$ & $-1$ \\
5  & 1 &  0 &  0 &  0 & $-1$ & $-1$ & 0 & 0 & 0 &  0 & 0 & 0 &  0 &  0 &  0 & 0 & 0 &  0 &  1 \\
11 & 1 & $-1$ & $-1$ & $-1$ & $-1$ & $-1$ & 1 & 1 & 1 &  1 & 1 & 1 &  1 &  1 &  1 & 1 & 1 &  1 &  1 \\
\bottomrule
\end{tabular}
\end{table}

Los coeficientes se estiman por mínimos cuadrados, $\hat{\bm{\beta}} = (\mathbf{X}^\top\mathbf{X})^{-1}\mathbf{X}^\top\mathbf{y}$, igual que en cualquier regresión lineal múltiple \citep{Montgomery2005}. En R el modelo se ajusta con la función \texttt{lm()}.

\subsubsection{Evaluación del modelo}

La significancia de cada término se evalúa con un análisis de varianza (ANOVA) con sumas de cuadrados tipo III, como en el artículo original. Para entender por qué se usa esta, conviene compararla con la suma de cuadrados tipo I, que es la que da la función \texttt{anova()} de R. La tipo I agrega los términos uno por uno, en el orden en que están escritos en la fórmula, y a cada término le asigna lo que reduce el error \emph{en el momento en que entra}. Si dos términos explican en parte lo mismo, el que entra primero se queda con esa parte y el que entra después recibe solo lo que agrega de nuevo; al cambiar el orden, cambian los resultados. Es como repartir el mérito de un trabajo en grupo según quién habló primero.

La tipo III evita ese problema. A cada término le asigna lo que reduce el error cuando se agrega a un modelo que ya tiene \emph{todos} los demás términos, es decir, solo lo que ese término explica y ningún otro puede explicar. Como cada término se compara siempre con el mismo conjunto (todos los demás), no importa el orden en que se escriban. Se calculó con la función \texttt{Anova()} del paquete \textsf{car} \citep{Fox2019}.

Además se usan las siguientes medidas:
\begin{itemize}
  \item \textbf{Prueba de falta de ajuste.} Las corridas repetidas en el punto central tienen exactamente las mismas condiciones, así que sus diferencias solo pueden deberse al ruido del experimento; esa variación se llama \emph{error puro}. La prueba compara el error del modelo con el error puro: si el modelo deja sin explicar mucho más que el ruido, su forma no es adecuada.
  \item \textbf{$R^2$ y $R^2$ ajustado.} La proporción de la variabilidad de la eficiencia que explica el modelo; el ajustado penaliza el número de términos.
  \item \textbf{$R^2$ predicho y PRESS.} El PRESS suma los errores que se cometen al predecir cada corrida con un modelo ajustado \emph{sin} esa corrida. Con él se calcula el $R^2$ predicho, que mide qué tan bien predice el modelo datos que no usó para ajustarse.
  \item \textbf{Desviación estándar ($s$) y coeficiente de variación.} El tamaño típico de los errores del modelo, en unidades de eficiencia y como porcentaje de la eficiencia media.
\end{itemize}

Los supuestos del modelo (errores normales, con varianza constante e independientes) se revisan con los residuales estudentizados. Un residual estudentizado es el residual dividido por una estimación de su desviación estándar, de modo que todos quedan en una escala comparable, parecida a la de una normal estándar. El internamente estudentizado usa la desviación estándar del modelo completo; el externamente estudentizado usa la de un modelo ajustado sin esa corrida, para que una corrida atípica no afecte la escala con la que se juzga a sí misma.
